\chapter*{Colophon}
\addcontentsline{toc}{chapter}{Colophon}

This book was produced from the \agentty{} repository using
\code{pdflatex} on Arch Linux. The body is set in Latin Modern Roman;
code listings use Latin Modern Mono; section headings use Latin Modern
Sans. Colours throughout come from a small palette --- \textcolor{mayablue}{blue},
\textcolor{mayalavender}{lavender}, \textcolor{mayagreen}{green},
\textcolor{mayarose}{rose}, \textcolor{mayaamber}{amber} --- carrying
consistent meaning across boxes and code highlights.

The four parts mirror four levels of zoom: foundations, architecture,
rendering, and reactivity. Each chapter is its own \code{.tex} file
under \file{chapters/}, included from \file{main.tex}. Re-ordering,
extending, or shortening the book is a matter of editing
\code{\textbackslash{}include} lines and re-running \code{make}.

The illustrative code snippets are simplified for exposition. Where
the simplification could matter --- the bit layout of a packed cell,
the exact arms of the \code{Cmd<Msg>} variant, the loop body of the
SIMD differ --- the prose points at the file and section in \maya{}
where the real version lives. The real version always wins.

\medskip

\textit{If a passage in this book disagrees with the source, the
source is right. Please file an issue against the \agentty{}
repository so we can fix the book.}

\medskip

\noindent Source layout for the book:

\begin{verbatim}
book/
  main.tex           -- root document, includes everything
  preamble.tex       -- packages, colours, code styles, boxes
  Makefile           -- `make` builds maya-book.pdf
  chapters/
    01-introduction.tex   -- Hello, maya
    02-terminal.tex       -- The terminal as a state machine
    03-cpp.tex            -- Modern C++ you need
    04-elm.tex            -- The Elm loop and algebraic effects
    05-dsl.tex            -- The compile-time DSL
    06-elements.tex       -- Elements, style, colour, border
    07-layout.tex         -- Flexbox in integers
    08-canvas.tex         -- The canvas and packed cells
    09-simd.tex           -- SIMD frame diffing
    10-ansi.tex           -- ANSI serialisation and the writer
    11-inline.tex         -- Inline mode and the witness chain
    12-signals.tex        -- Signals, computeds, effects
    13-events.tex         -- Events, focus, input parsing
    14-cmdsub.tex         -- Cmd, Sub, and the runtime
    15-widgets.tex        -- The widget catalogue
    16-build.tex          -- Building, packaging, shipping
    99-colophon.tex       -- this page
\end{verbatim}

\medskip

\noindent To rebuild from scratch:

\begin{verbatim}
cd book/
make             # runs pdflatex twice, regenerates maya-book.pdf
\end{verbatim}

\vfill

\begin{center}
\small\itshape
Written for anyone who has ever opened a terminal and wondered
what was on the other end of the wire.
\end{center}
